%! @@TITLE@@ — Unit @@UNITINT@@, Lesson @@LESSONID@@ — @@DOCTITLE@@
\documentclass[10pt]{article}
\usepackage{@@PREFIX@@-article}
\usepackage{@@PREFIX@@-boxes}
\newcommand{\TallMath}[1]{$\displaystyle #1\rule[-1.4em]{0pt}{3.2em}$}

\begin{document}

\pageheader{Unit @@UNITINT@@, Lesson @@LESSONID@@}{@@DOCTITLE@@}
% NAMESTRIP: no \namedateperiod here. The student writes their name once, on the
% cover this component is stapled behind. See LESSON_SHAPE.md (namestrip).
\vspace{0.15in}

% TODO: author the @@DOCTITLE@@. See templates/lesson/components.md for the expected structure.
% PILOT SHAPE (2026-09-06) — mirror unit01/lesson02. GUIDED NOTES & PRACTICE carry the whole
% gradual release and END AT THE GUIDED PRACTICE:
%   objectivebox (printed targets, worded as on the cover) -> vocabbox of \termblank{} terms
%   (key: \vocabans{}{} with a TWO-LINE definition so the boxes match in height) ->
%   EXACTLY TWO \begin{notesbox}{N. Title} sections the teacher models (I DO), each in two
%   moves, the second move introduced by \textbf{\textcolor{cerulean}{Its Title.}} ->
%   one \begin{practicebox} ("Guided Practice", WE DO) of three or four lettered parts, the last
%   testing the misconception on new ground.
% No hookbox (the hook is a slide), no solo practice box, no extensionbox. The crux is the second
% half of section 2. Target 3-4 pages, blank and key. The homework is the individual practice.
% HOMEWORK: a remindbox ("This is your graded homework..."), a scenariobox context, ~6 items in a
% notesbox (split into a "..., continued" box where the pages fall — never strand a lead-in from
% its table), a contrast pair, the crux head-on, one spiral item, and a closing \begin{spiralbox}
% preview. Scored; started in class in the last eight minutes.
% GENERATE THE KEY from the blank: python3 templates/lesson/mkkey.py blank key answers.json
% (see templates/lesson/examples/). Never hand-edit a key.
% \boxguard before every breakable box — and in the _key too. Raise it (\boxguard[30]) when
% the box opens with an unbreakable TikZ figure or tabularx.
% Multi-step solutions go in a \begin{work} block authored BYTE-IDENTICALLY here and in the
% _key: the blank reserves its exact height, the key prints it. That is the work rule, and it
% is what keeps this component the same number of pages as its key.
\boxguard
\begin{notesbox}{@@DOCTITLE@@}
\writelines{4}
\end{notesbox}

\end{document}
